\section{卷积把平移结构写进线性映射}
\label{sec:14-Deep-Learning/04-Convolutional-Networks/01}

拿一张 \(224\times 224\) 的 RGB 图像。不是什么大图——手机随手拍的尺寸。三个通道，摊平是 \(224\times 224\times 3 = 150{,}528\) 个数字。接一层只有 64 个输出神经元的全连接层，权重矩阵 \(64\times 150{,}528\)，大约 963 万个参数。

963 万个参数，每一个都是独立的旋钮。左上角第 \((3,7)\) 个像素连到第 15 个输出神经元的权重，和右下角第 \((200,150)\) 个像素连到第 15 个输出神经元的权重，在数学上是两个互不相干的自由变量。训练数据里如果两个位置都出现过竖边，网络必须分别为它们各自学会竖边长什么样——学两遍，互不借力。位置隔了两万个像素，参数就隔了两万行。而 \(224\times 224\) 的图里有将近五万个空间位置，每个位置都要从头来过。

你可能会想：数据够多不就行了？把输入分辨率翻倍到 \(448\times 448\)。输入维度变成约 60 万，同一层全连接的参数从 963 万暴涨到约 3840 万。再翻倍，冲向 1.5 亿。这还只是一层。参数数量的平方增长不只是一个显存问题——它在问一个更根本的东西：图像的``空间位置''真的需要被当作独立变量来建模吗？

一条竖边无论出现在左上角还是右下角，都是同一种竖边。局部邻域的像素相对排列没变，只是绝对坐标变了。如果网络在某个位置学会了识别竖边，把同一组权重搬到另一个位置，它应该也能识别。问题是怎么让这件事发生。

靠训练数据来``希望''网络自己发现位置间的参数可以共享，取决于优化路径、初始化和数据分布的偶然。不同位置的权重完全可能分道扬镳——没有结构性的力量把它们绑在一起。

出路不是在训练中祈祷参数收敛到相似的数值，而是把``同一组参数被所有位置使用''直接写进算子本身。让函数类从起点就只能表达那些具有平移重复结构的映射。这就是卷积。

\subsection{同一组权重核被所有位置共享}

设输入序列 \(x_0,\ldots,x_{T-1}\)，核 \(w_0,\ldots,w_{K-1}\)。深度学习库里的``卷积''绝大多数实现为互相关：

\[
y_t = \sum_{j=0}^{K-1} w_j\, x_{t+j}.
\]

数学上的卷积定义会翻转核索引：\(\sum_j w_j x_{t-j}\)。但学习任务不需要区分这两种定义——学出来的核只要把参数顺序反过来就能相互转换，梯度下降自己会找到合适的排列。工程中统称为卷积。真正要紧的是：同一组 \(w_0,\ldots,w_{K-1}\) 被每一个输出位置 \(t\) 共享。

拿一个具体的一维信号走一遍。\(x = [2,\,1,\,4,\,3,\,0,\,5]\)，核 \(w = [1,\,0,\,-1]\)——一个差分检测器：当前位置减两步后位置的值。位置 \(t=0\)：\(y_0 = 1\cdot 2 + 0\cdot 1 + (-1)\cdot 4 = -2\)。位置 \(t=3\)：\(y_3 = 1\cdot 3 + 0\cdot 0 + (-1)\cdot 5 = -2\)。同一个核，在两个不同的局部窗口里，各自算出了差分。不需要为位置 0 和位置 3 分别学一套参数。

核在输入序列上滑动的姿态由三个参数控制。步幅 \(s\)：每次右移多少格，\(s=2\) 意味着隔一个位置算一次，输出长度大致减半。padding \(p\)：在序列两端各补 \(p\) 个位置（通常是零），让边缘附近不至于被截断得太快。dilation \(d\)：在核元素之间插入空位。\(d=2\)、\(K=3\) 时，三个参数的实际覆盖范围是 \(2\times 2+1=5\) 个输入位置——不增加参数数量，用稀疏采样换取更大的空间视野。

引入 dilation 后，有效核宽度 \(K_{\mathrm{eff}} = d(K-1)+1\)。输出长度公式：

\[
T_{\mathrm{out}} = \left\lfloor\frac{T + 2p - K_{\mathrm{eff}}}{s}\right\rfloor + 1
\]

不需要背。想象在长度为 \(T+2p\) 的数组上滑动宽度 \(K_{\mathrm{eff}}\) 的窗口，每次右移 \(s\)。第一个窗口起始位置是 0，最后一个窗口起始位置不能超过 \((T+2p)-K_{\mathrm{eff}}\)，否则右端超出数组。从 0 开始，每次跳 \(s\)，能完整放下的窗口个数就是 floor 部分加 1。若除法不能整除，floor 把零头砍掉——数组末尾有几个位置永远不会成为完整窗口的起点，它们的信息在这一层被丢弃了。

\subsection{通道维是张量缩并的另一个轴}

推广到二维。对 \(C_{\mathrm{in}}\) 个通道、高 \(H\)、宽 \(W\) 的输入 \(X\in\R^{C_{\mathrm{in}}\times H\times W}\)，卷积核是四维张量 \(W\in\R^{C_{\mathrm{out}}\times C_{\mathrm{in}}\times K_h\times K_w}\)。忽略 stride 和 padding：

\[
Y_{c,i,j} = b_c + \sum_{r=1}^{C_{\mathrm{in}}}\sum_{u=0}^{K_h-1}\sum_{v=0}^{K_w-1} W_{c,r,u,v}\, X_{r,\,i+u,\,j+v}.
\]

从张量缩并的视角：空间窗口索引 \(u,v\) 和输入通道 \(r\) 在求和后消失（被缩并），输出通道 \(c\) 和空间位置 \(i,j\) 保留为新张量的轴。一个 \(3\times3\times 3\) 的核（\(C_{\mathrm{in}}=3\)，RGB），在每个空间位置与输入的 \(3\times K_h\times K_w\) 小块做内积，吐出一个标量。64 个这样的核产生 64 个输出通道——每个核从同一个输入邻域里提取不同的特征模式：一个对竖边敏感，另一个对横边敏感，第三个对红-绿对比敏感。

参数总量：\(C_{\mathrm{out}}C_{\mathrm{in}}K_hK_w + C_{\mathrm{out}}\)（加的是 bias）。与 \(H,W\) 完全无关。一张 \(1000\times 1000\) 的图像和一张 \(32\times 32\) 的图像，同一层卷积的参数数量一模一样——差别只在计算量随空间尺寸线性增长。而若把输入摊平接全连接，\(H\) 或 \(W\) 翻倍，参数就要平方级增长。参数共享同时降低了计算量和函数类容量：计算便宜和表达受限是同一枚硬币的两面。

\subsection{卷积就是结构极度规整的矩阵乘法}

把卷积写成线性映射 \(y = C x\) 可以直观看到参数共享到底长什么样。一维，\(T=4\) 个输入，\(K=3\) 的核，stride 1、无 padding——输出长度 \(4-3+1=2\)。矩阵 \(C\) 为：

\[
C = \begin{bmatrix}
w_0 & w_1 & w_2 & 0 \\
0 & w_0 & w_1 & w_2
\end{bmatrix}.
\]

每一行是同一组权重向右移动一个位置——这是一个 Toeplitz 矩阵：每条对角线上的元素相等。\(w_0\) 出现了 2 次，\(w_1\) 和 \(w_2\) 也各出现了 2 次。全连接层的矩阵不会有任何这种重复约束。二维卷积对应的矩阵是 doubly block Toeplitz，每组 \(C_{\mathrm{out}}\times C_{\mathrm{in}}\) 个空间核各自对应一块这样的结构。

从这个矩阵形式能读出 padding 的效果：加一个零填充，矩阵形状改变，边界行的内容也变了——手工填的零进入了原本应该看到真实像素的位置。

\subsection{平移等变性是算子自带的代数性质}

记平移算子 \((T_a x)_t = x_{t-a}\)。在无限序列或忽略边界效应时，stride 为 1 的卷积算子 \(C_w\) 满足：

\[
C_w(T_a x) = T_a(C_w x).
\]

验证只需展开右侧第 \(t\) 个分量：\(\sum_j w_j x_{t+j-a} =\) 左侧第 \(t-a\) 个分量。输入平移 \(a\) 个位置，输出跟着平移 \(a\) 个位置——输入和输出之间的相对空间关系保持不变。这叫等变，不是不变。输出变了，但和输入变得一样多。

等变性是算子自身的代数性质——用随机初始化的权重、没经过任何训练也成立。它不是从数据中学来的。

若最终层接一个全局平均池化，把整个空间维压成一个数，输入平移导致的输出平移被均值抹掉了——输出不再随输入平移而改变。这才叫（近似）不变。分类需要这种最终的不变性，但语义分割和目标检测依赖空间信息不被过早销毁，中间层更依赖等变。

这条等式并非全局成立。零填充在边缘制造了不来自真实数据的输入——不是学出来的特征，而是硬塞进去的常数 0。取 \(x=[1,2,3]\)，核 \(w=[1,1]\)，不填充时输出 \([3,5]\)。向左平移 1 步后，在有限数组上没有定义，必须借助填充。使用零填充后，边界位置的窗口看到了手工填入的零，原本不该出现的信息混了进来。stride 大于 1 时，平移一个像素可能导致输出落在不同的下采样网格上——相邻的两个输入位置，一个被 stride 捕获，另一个被跳过，在输出表示里的命运截然不同。池化和有限裁剪进一步累积这些误差。实际 CNN 只在相应的离散平移和边界条件范围内近似等变。

\subsection{感受野用深度换取视野}

一层 \(3\times3\) 卷积，每个输出像素只看输入的 \(3\times3\) 邻域。堆两层：第二层的每个像素看第一层的 \(3\times3\)，而第一层这 9 个像素每个各自看了输入的 \(3\times3\)——间接覆盖 \(5\times5\)。再堆一层变成 \(7\times7\)。递推公式：

\[
j_\ell = j_{\ell-1}s_\ell,\qquad
r_\ell = r_{\ell-1} + (K_\ell-1)j_{\ell-1},
\]

初始 \(j_0=r_0=1\)。stride 全为 1 的 \(3\times3\) 堆叠：\(r_1=3\)，\(r_2=5\)，\(r_3=7\)，\(r_4=9\)，\(r_5=11\)。五层拥有 \(11\times11\) 的感受野。

但和一层 \(11\times11\) 卷积相比：五层 \(3\times3\) 参数总量 \(5\times C_{\mathrm{out}}C_{\mathrm{in}}\times 9\)，而后者是 \(C_{\mathrm{out}}C_{\mathrm{in}}\times 121\)——参数多出约 2.7 倍，中间还少了四层非线性。堆小核不仅省参数，还在相邻像素之间插入了可学习的非线性决策。

理论感受野只标出``存在一条从输入像素到输出单元的路径''，不保证这条路径实际承载有效梯度。训练后中心像素附近的权重通常远大于边缘——有效感受野远小于理论值，呈近似高斯衰减。这条递推、空洞卷积对它的指数级放大，以及有效宽度为何只按层数的平方根增长，下一篇逐条说清。

\subsection{分组卷积和 depthwise 把通道耦合拆分}

标准卷积中，每个输出通道与所有输入通道全连接。分组卷积把输入通道分成 \(g\) 组，每组独立计算——组内全连接，组间完全隔离。参数从 \(C_{\mathrm{out}}C_{\mathrm{in}}K_hK_w\) 降到 \(C_{\mathrm{out}}C_{\mathrm{in}}K_hK_w/g\)。极端情况 \(g=C_{\mathrm{in}}=C_{\mathrm{out}}\)：每个通道单独做空间卷积，通道之间零交互——这叫 depthwise convolution。

后面再接 \(1\times1\) 卷积（pointwise），核尺寸 \(1\times1\) 不做空间混合，纯粹在通道间做线性组合。这种 depthwise separable convolution 把空间混合和通道混合拆成两步独立运算。总计算量从 \(C_{\mathrm{out}}C_{\mathrm{in}}K_hK_wHW\) 降到约 \(C_{\mathrm{in}}K_hK_wHW + C_{\mathrm{out}}C_{\mathrm{in}}HW\)。对 \(K_h=K_w=3\) 和较大的 \(C_{\mathrm{out}}\)，加速比接近 9 倍。MobileNet 利用这种分解在手机端跑了接近当时大模型的精度。

让出的是单层内部通道间不再直接交换信息。标准卷积中输出通道 1 能同时看到输入通道 1 的边缘信号和输入通道 2 的颜色信号并做融合；depthwise separable 里，depthwise 那步输出通道 1 只能看到输入通道 1 的空间模式，通道融合被推迟到随后的 pointwise 层——如果 pointwise 层本身又很窄，融合能力就受限于这个窄瓶颈。计算便宜不等于结构上没有损失。

\subsection{平移共享的假设在什么地方会撞墙}

卷积把``相同局部模式在不同位置具有相近含义''写进了算子。自然界图像近似满足：竖边就是竖边。但以下场景不会配合这个假设：

\begin{itemize}
\tightlist
\item
  表格数据。列之间没有稳定的邻接顺序。第三列和第四列的``局部邻域''与第五列和第六列的``局部邻域''语义上没有共享的理由——除非列的排列编码了物理拓扑。按任意列序做卷积等于碰运气。
\item
  绝对位置决定语义。棋盘坐标 (e4) 和 (a1) 意义完全不同；UI 界面中左上角的关闭按钮和右下角的确认按钮；CT 片中器官的固定解剖位置——这些场景中，同一个模式出现在不同位置可能意味着不同事物。强制权重共享是在迫使网络花额外的非线性层去编码位置差异——用容量弥补一个错误的归纳偏置。
\item
  旋转和缩放。倾斜 45 度的边缘对水平卷积核是全新模式，需要从零学起。旋转不在平移群内，普通卷积的等变性不覆盖旋转。数据增强是退而求其次的弥补。
\item
  部署域漂移。训练数据中物体总是居中，推理时偏移到角落。边界附近填充伪影与训练分布不同——卷积的等变性只保证训练分布内的平移关系，不保证域外有效。角落识别精度可能系统性低于中心。
\end{itemize}

架构设计把假设直接写进了线性算子。数据只能在假设划定的范围内修正参数。如果假设本身与问题的物理结构冲突，堆更多层、加更多数据，不过是在逼近一个被错误约束的函数类。开始训练之前先确认：位置共享是否与任务的物理结构一致、边界填充是否在制造训练时不曾出现的新特征、下采样模式是否兼容小尺寸目标、训练分布中的平移规律在部署条件下是否仍然成立。

接下来的四篇把这份清单逐项展开：先说清视野能推多远，再看下采样与上采样各自赔进去什么，最后回头给这整套先验估一个价。
